% sections/01_intro.tex -- Section 1, Introduction (W4). Plan: BLUEPRINT 2.1 (argument, results cells, four literature paragraphs, "do not say" list),
% 5.3 (financing framing), 9 (register), Overrides O4 (Trustees figure verified by web search on 2026-11-09), O5 (provenance comments), O8 (5.5-page budget).
% No equations. "Rigid" and "flexible" appear once each, in 1.1, as plain adjectives (BLUEPRINT Decision 2). Every number carries a "% src:" comment.
% Abbreviations: T: = paper/tables/; 1:NN, 5:NN, 7:NN, 11:NN, 12:NN = frot/sub2025/{1_intro,5_quantitative_model,7_policy_reforms,11_app_data,12_app_quant}.tex;
% v6: = frot/current/.
\section{Introduction}\label{sec:intro}

\subsection{The financing problem and two ways to change the claiming schedule}\label{subsec:financing}

The Old-Age, Survivors, and Disability Insurance (OASDI) program is not financed over the long run at its scheduled benefits and tax rates: the 2024 Trustees Report puts its 75-year actuarial deficit at 3.50 percent of taxable payroll \citep{trustees2024}. % src: v6:6_quantitative_wedges_resubmission_v6.tex:546; verified against the 2024 report by web search (Override O4)
The response most often legislated or proposed is a higher retirement age. The 1983 amendments raise the normal retirement age (NRA), the age at which a retiree receives the full benefit, from 65 to 67 for cohorts born in 1960 and later, and the Congressional Budget Office has analyzed an option that raises it to 70 for cohorts born in 1981 and later \citep{cbo2024socialsecurity}. % src: 1:44 (footnote)

A higher NRA is one of two ways to change the schedule that relates a retiree's benefit to the age at which it is first claimed. Under the rules faced by the cohorts studied here, benefits can be claimed from age 62; the benefit falls by 6.67 percentage points of the full amount for each year of claiming before the NRA of 65 and rises by a delayed retirement credit (DRC) of 4.5 percentage points for each year of claiming after it, up to age 70. % src: 7:90; T:SS_DRC.tex:10; T:parameters_ss.tex:46-48
Raising the NRA shifts this schedule: the benefit at every claiming age falls, whether or not anyone changes the age at which they claim. The second way leaves the benefit at the NRA unchanged and changes the slope of the schedule around it, through a higher DRC or a steeper adjustment on both sides of the NRA---a steeper actuarial adjustment of the benefit to the claiming age. Benefits then change only for individuals who claim away from the NRA: those who delay receive more, and those who claim early receive the same or less. The first instrument makes the retirement age more rigid; the second rewards a flexible claiming age without moving the statutory one.

If no one changed behavior, a higher NRA would save money, because every benefit falls; a higher DRC would cost money, because every late claimant receives more; and a symmetric adjustment would save or cost depending on how many claim before and how many after the NRA. Which reform saves more per unit of welfare lost, or whether one of them saves money while raising welfare, depends on who changes their claiming age and their labor supply in response, and on what happens to the taxes they pay and the other benefits they receive. That is an empirical question about the behavior of a population that differs in health, wages, assets, and expected longevity, and this paper answers it in an estimated life-cycle model.
% src: T:DRC_visualization_4.tex:37-58 (AN factors lower at every age; AD factors higher after the NRA only; AL factors lower before and higher after the NRA); sign inferences of \autoref{sec:suffstats}

\subsection{What we do}\label{subsec:whatwedo}

We estimate a life-cycle model of labor force participation, hours, saving, and Social Security claiming in which individuals face health, mortality, wage, and medical-expense risk, in the tradition of \citet{french2005effects} and \citet{french2011effects}. Health takes three states, good, poor, and disability; it governs mortality, out-of-pocket medical spending, and wages, and it changes the fixed time cost of participation, which is estimated separately for healthy and for unhealthy workers. Exit from the labor force is reversible at an estimated reentry cost, and claiming is irreversible. The model includes the claiming-age schedule, the benefit formula, the retirement earnings test, disability insurance, Medicare, defined-benefit pensions, a progressive income tax, the payroll tax, and a consumption floor. It is estimated in two steps on the 1920--1940 birth cohorts of the Panel Study of Income Dynamics (PSID), with medical spending from the Medical Expenditure Panel Survey (MEPS). The first step estimates health transitions, survival, medical spending, and wage profiles by health; the second estimates seven preference parameters by the method of simulated moments on life-cycle profiles of participation and hours by health, assets, and reentry. Health reaches the reform results through the first-step profiles---mortality, medical spending, and wages by health---and through the health-dependent participation cost; the reform tables report participation by age group, and \autoref{subsec:participation} states what the profiles imply, in the model, about who moves. % src: 11:18; 5:533; T:estimated_parameters.tex:19-21 (participation cost by health); BLUEPRINT Decision 3 (health reading stated once, in 7.3)

We then use the model to evaluate seven families of reforms with one welfare measure and one fiscal accounting. Four change Social Security: a higher NRA; a higher DRC; a symmetric steepening of the schedule around the NRA; and increases in the three bracket rates of the benefit formula, one at a time or together. Three change other programs: expansions of Medicare coverage after age 65; reductions in income taxes after age 62; and subsidies to retirement saving, read as changes in the tax treatment that retirement saving already receives through individual retirement accounts and 401(k) plans. % src: seven families and their definitions: T:welfare_age_dependence_mvpfv5.tex:108-126, T:welfare_progressivity_mvpfv7.tex:115-158, T:welfare_medicare_mvpfv8.tex:113,116, T:welfare_tax_mvpf_v8.tex:114-118, T:welfare_savings_mvpf_v3.tex:120-123; ages 62 and 65: T:parameters_ss.tex:46-47; 7:423 (IRAs and 401(k) plans); BLUEPRINT Decision 9 (no count of experiments)
Welfare is the consumption-equivalent change in the ex ante expected utility of the cohort at age 30. The fiscal accounting separates the change in the outlays of the program being reformed from the change in every other tax and transfer the cohort pays or receives, and combines them into the change in the government's net spending on the cohort, all as per-capita present values at age 62. The marginal value of public funds (MVPF) of each reform is the ratio of the cohort's willingness to pay to the change in net government spending \citep{hendren2020unified}. Because the model solves the individual's problem under each set of rules, the willingness to pay is exact for discrete reforms with fixed participation costs and irreversible claiming rather than a first-order approximation, and the fiscal response of every program is computed jointly with it. \autoref{sec:accounting} maps each column of the reform tables into the MVPF and relates the structural MVPF to the consumption-smoothing benefit and the fiscal multiplier of the sufficient-statistics literature.
% src: 5:23 (entry age 30); 12:611 (NPVs at 62)

The exercise takes the perspective of one cohort. Reforms are permanent changes in the rules it faces from age 30; fiscal savings are valued at face value, with a unit welfare weight on their recipients outside the cohort; no transition is modeled; and the MVPF is reported as a shadow price to which any intergenerational weight can be applied. A reform contributes to closing the Trustees' gap if it lowers the government's net spending on the cohort.
% src: 5:23 (entry age 30); 12:611 (NPVs at 62); BLUEPRINT 5.3 (incidence stance)

\subsection{Results}\label{subsec:results}

Raising the NRA to 66, 67, or 68 lowers the cohort's willingness to pay by 0.38, 0.74, and 1.09 percent of consumption, lowers the present value of its retirement benefits by 6.13, 11.82, and 16.75 percent, and lowers net government spending by 0.87, 1.75, and 2.38 percent. % src: T:welfare_age_dependence_mvpfv5.tex:108-110
The ratio of the welfare loss to the fiscal saving is the MVPF of reversing the reform: lowering the NRA from 66, 67, or 68 back to 65 would deliver 0.17, 0.16, or 0.18 dollars of welfare per dollar of net spending. % src: T:welfare_age_dependence_mvpfv5.tex:108-110; 7:119 (reversal reading)
Raising the DRC from the cohort's 4.5 percent to 8--10 percent per year raises the willingness to pay by 0.16 to 0.3 percent of consumption and lowers retirement-benefit outlays by 2.46 to 5.05 percent; the government's other net outlays fall as well, by 0.04 to 0.10 percent, and net spending falls by 0.44 to 0.78 percent. % src: T:welfare_age_dependence_mvpfv5.tex:112-116 (Net -0.45, -0.44, -0.66, -0.73, -0.78); T:DRC_visualization_4.tex:43-47
Setting equal early-claiming penalties and delayed credits of 7 to 15 percent per year raises the willingness to pay by 0.10 to 0.80 percent of consumption and lowers net spending by 0.12 to 1.57 percent. % src: T:welfare_age_dependence_mvpfv5.tex:118-126
These reforms raise ex ante welfare and generate a fiscal surplus, so their MVPF is infinite, and participation at ages 65--69 rises under every one of them, by 0.22 to 2.61 percentage points. % src: T:welfare_age_dependence_mvpfv5.tex:112-126
Two pairs hold the fiscal saving approximately fixed. An NRA of 66 and a symmetric adjustment of 10 percent per year lower net spending by 0.87 and 0.86 percent; the willingness to pay is $-0.38$ percent of consumption under the first and $+0.29$ under the second. % src: T:welfare_age_dependence_mvpfv5.tex:108,121
An NRA of 67 and a symmetric adjustment of 15 percent lower it by 1.75 and 1.57 percent, with willingness to pay of $-0.74$ and $+0.80$. % src: T:welfare_age_dependence_mvpfv5.tex:109,126
A benefit increase for late claimants lowers net spending; a benefit cut for every claimant lowers it by a similar amount and lowers welfare. Participation responses by age locate the difference. Under a steeper adjustment participation rises after the NRA under every reform and, beyond the smallest steps, falls at ages 62--64: the reform moves work and claiming from before the NRA to after it, and the steepest adjustments shift the incidence toward those who claim early. % src: T:welfare_age_dependence_mvpfv5.tex:112-126 (62-64 negative from AD3 and AL3 on; 65-69 positive in every row)

The quantitative analysis opens by removing each policy instrument in turn. Eliminating Social Security benefits lowers welfare by 5.89 percent of consumption; the loss from incomplete Medicare coverage after age 65 is 4.28 percent; raising the net return on saving to the gross return raises welfare by 3.60 percent and eliminating the income tax by 3.16 percent; eliminating disability benefits lowers welfare by 0.34 percent. % src: T:welfare_decomposition_gradient.tex:27,28,33,36,37
Social Security, the instrument with the most welfare at stake, is taken up first.

Within reforms that raise the generosity of retirement benefits, the first bracket of the benefit formula is where a dollar of net spending buys the most welfare. Increases in the first bracket rate have MVPFs from a lower bound of 0.52, for the smallest step, which is approximately self-financing, to 1.14; proportional increases in all three rates have MVPFs of 0.30 to 0.61, and increases in the third bracket 0.29 to 0.30. % src: T:welfare_progressivity_mvpfv7.tex:115,120; :149-158; :137-146
Expansions of Medicare coverage after age 65 by 20 percentage points, for everyone or for nonparticipants only, have MVPFs of 0.146 and 0.155. % src: T:welfare_medicare_mvpfv8.tex:113,116
Reductions in income taxes after age 62 raise participation at every age and have MVPFs of 0.23 to 0.26. % src: T:welfare_tax_mvpf_v8.tex:114-118
Subsidies to retirement saving of 0.2, 0.4, and 0.8 percentage points of return have MVPFs whose lower bounds, computed against the subsidy alone, are 0.96, 0.92, and 0.84. % src: T:welfare_savings_mvpf_v3.tex:120,122,123
Across the menu, the reforms that lower net spending are the claiming-schedule reforms and the smallest first-bracket increase; every other reform raises welfare and net spending together, and the MVPF orders them. % src: T:welfare_age_dependence_mvpfv5.tex:108-126; T:welfare_progressivity_mvpfv7.tex:115-158; T:welfare_medicare_mvpfv8.tex:113,116; T:welfare_tax_mvpf_v8.tex:114-118; T:welfare_savings_mvpf_v3.tex:120-123 (signs of the Net column)

\subsection{Relation to the literature}\label{subsec:literature}

The model belongs to the class of estimated life-cycle models of retirement with health and medical-expense risk. \citet{Rust1997} study Social Security and Medicare under incomplete markets; \citet{french2005effects} estimates the effects of health, wealth, and wages on labor supply and retirement; \citet{french2011effects} add employer-provided and public health insurance. Other members of the class study retiree health insurance \citep{Blau2008b}, Social Security and the saving of low-income households \citep{vanderklaauw2008social}, medical expenses and the lifetime costs of bad health \citep{Denardi2010,Denardi2018}, disability insurance \citep{Low2015}, and medical spending and the evolution of health over the life cycle \citep{Pashchenko2019,hosseini2022evolution}. \citet{Haan2014}, \citet{laun2019health}, \citet{Li2018}, \citet{iskhakov2021effects}, \citet{Jones2023}, \citet{Bairoliya2023}, \citet{Borella2023}, and \citet{Yu2024} use models of this class to evaluate pension and Social Security reforms, among them increases in the retirement age and changes in the claiming-age adjustment and the benefit formula; \citet{blundell2016retirement} survey the literature. Relative to \citet{french2005effects} and \citet{french2011effects}, the model here lets the fixed cost of participation depend on health, models disability insurance explicitly, allows work at every age until death, and is estimated on the general population rather than on male household heads (\autoref{subsec:scope}).

Social Security reform has been studied in general-equilibrium overlapping-generations models \citep{nishiyama2007does,Imrohoroglu2012,kitao2014sustainable} and in life-cycle models of labor supply and social insurance \citep{huggett2010well,Erosa2012}; \citet{Imrohoroglu2012} and \citet{kitao2014sustainable} include an increase in the normal retirement age among the reforms they evaluate. An empirical literature studies the US and Swiss increases in the statutory age \citep{mastrobuoni2009labor,behagel2012,deshpande2024sticky,lalive2023social}, the response of retirement to financial incentives and future entitlements \citep{coile2007future,manoli2016nonparametric}, the gain from delaying claiming and the means to finance a delay \citep{shoven2014does,goda2018financial}, how claiming decisions are made and framed \citep{shoven2017,brown2016framing,liebman2015would,seibold2021reference}, and old-age support and the earnings test \citep{fetter2018government,gelber2020employment}. \citet{gruber1998social} and the volumes they edited \citeyearpar{gruber1999social,Gruberwise2002,gruber2009social} compare countries; \citet{Coile2019} documents the US trend toward working longer in a later volume of the same series. This paper evaluates the two instruments for one cohort, at fixed prices, with the same welfare measure and the same fiscal accounting for both, and reports who moves by age group.

The accounting draws on the marginal value of public funds of \citet{hendren2016policy}, \citet{hendren2020unified}, and \citet{finkelstein2020welfare}, whose antecedents are the excess-burden and marginal-cost-of-funds measures of \citet{mayshar1990measures}, \citet{slemrod1996costs}, and \citet{kleven2006marginal}, and which \citet{finkelstein2008did} and \citet{finkelstein2019subsidizing} apply to Medicare and to health-insurance subsidies. The sufficient-statistics approach to social insurance \citep{baily1978some,chetty2006general,chetty2008moral,chetty2009sufficient,chetty2013social,kleven2021sufficient} values a program by the consumption-smoothing benefit of its transfers against their mechanical and behavioral fiscal cost; \citet{landais2021value} and \citet{kolsrud2024retirement} apply it to unemployment insurance and to pension design, and \citet{haller2022welfare} to pension reforms, summarizing the cost side by a fiscal multiplier. \autoref{sec:accounting} writes the structural MVPF as the ratio of one plus a consumption-smoothing term to a fiscal multiplier and states which of these objects the reform tables identify.

The companion paper \citep{ndiaye2026flexible} characterizes the constrained-efficient allocation in the same economy and decomposes retirement behavior into the contributions of participation costs, health-related productivity differences, and mortality, in the dynamic optimal-taxation tradition of \citet{diamond1978model}, \citet{golosov2006designing}, \citet{farhi2013insurance}, \citet{golosov2016redistribution}, \citet{Stantcheva2017}, and \citet{shourideh2015incentives}, and of the work on retirement and taxation of \citet{cremer2004social}, \citet{grochulski2010nonseparable}, \citet{michau2014optimal}, \citet{chone2014income}, and \citet{hosseini2019retirement}. This paper uses the estimated model as a laboratory for reforms of the statutory rules; it does not use the constrained-efficient allocation as a benchmark, and it makes no claim that any reform in the menu approximates it.

\subsection{Scope and roadmap}\label{subsec:roadmap}

The analysis is partial equilibrium for one cohort at fixed factor prices. \autoref{subsec:scope} lists the further restrictions---an individual decision-maker, accidental bequests, health observed by the government, time-separable preferences, and statutory ages that act only through incentives---and \autoref{sec:scope} states how each bears on the results and gives a reduced-form expression for how factor-price responses would enter the fiscal accounting.

\autoref{sec:institutions} describes the rules faced by the 1920--1940 cohorts, the two ways of changing the claiming schedule, the reform menu, and the cohort perspective.
% src: 11:18 (1920--1940 birth cohorts)
\autoref{sec:model} presents the model and \autoref{sec:estimation} the data, the estimation, and the fit, with the differences by health on which the reforms act. \autoref{sec:accounting} defines the welfare measure, the fiscal columns, and the MVPF, and relates them to the sufficient-statistics objects. \autoref{sec:elimination} removes each instrument in turn. \autoref{sec:claiming} evaluates the claiming-schedule reforms, \autoref{sec:progressivity} the benefit-formula reforms, and \autoref{sec:other} the Medicare, income-tax, and saving-subsidy reforms. \autoref{sec:menu} places every reform side by side. \autoref{sec:scope} discusses scope and sensitivity, and \autoref{sec:conclusion} concludes. \autoref{app:data} to \autoref{app:derivations} contain the data and first-step estimation, the solution method and the fit of assets and reentry, the construction of the reforms, and the derivations.
